\section{Introduction and statements}
Character-twisted moments of the Legendre family give concrete finite sums
whose cohomological interpretation depends on boundary terms and arithmetic
normalizations. This paper treats cubic characters and the second and fourth
moments. Its central purpose is an explicit construction of the associated
representations and their Frobenius formulas, including the inert-prime sign. The consolidated local arguments now
also determine the two bad-prime factors and conductor exponents.
The exact finite-sum prescription underlying the numerical degree-four
recipe in the earlier preprint \cite{WalkeV1} is thereby identified
at good primes with a specified representation.

\subsection{Notation and Frobenius convention}
Put $K=\Q(\zeta_3)$, $E=\Q_2(\zeta_3)$ and
$U=\PP^1\setminus\{0,1,\infty\}$. Let
$\pi:\mathcal E\to U$ be the smooth projective Legendre family
\[
 \mathcal E_\lambda:y^2=t(t-1)(t-\lambda),\qquad
 \mathcal V=R^1\pi_*E.
\]
Let $\mathcal L_\rho$ be either nontrivial character summand of the actual
arithmetic cover $z^3=\lambda$. No additional constant Galois twist is
allowed in this definition. For an integer $r\in\{2,4\}$ set
\[
 H_r=H^1_!(U_{\overline K},\Sym^r\mathcal V\otimes\mathcal L_\rho),
 \qquad H^1_!=\operatorname{im}(H_c^1\longrightarrow H^1),
 \qquad W_2=H_2,\quad V_4=H_4.
\]
Frobenius always means \emph{geometric Frobenius on actual cohomology}.
Write $\tau$ for the character of $\Q_2(-1)$, so
$\tau(\Frob_{\mathfrak p})=N\mathfrak p$ and
$\tau=\chi_{\rm cyc}^{-1}$ in the usual convention on $\Q_2(1)$.
The modular comparison uses the cohomological dual of the conventional
arithmetic-Frobenius Deligne representation.

For a good residue field of cardinality $q$ and $\lambda\ne0,1$, put
\[
 a_q(\lambda)=q+1-\#\mathcal E_\lambda(\F_q),\qquad
 M_{3,j}(q)=\sum_{\lambda\in\F_q\setminus\{0,1\}}
                  \rho(\lambda)a_q(\lambda)^j.
\]
Here $\rho$ is the trace character of the specified Kummer summand.
When taking powers of one Frobenius, lift it by the finite-field norm.
At an inert rational prime of $K$, $\F_{p^2}$ is the initial residue
field; no nontrivial cubic character of $\F_p^\times$ is being lifted.
Define
\[
 v(q)=-M_{3,4}(q)+3qM_{3,2}(q)+q^2-1.
\]

\begin{theorem}[Summary of the fourth-moment results]\label{thm:main}
The representation $V_4$ has dimension four over $E$ and is unramified
outside the primes over $2,3$. At every good prime of $K$,
$\Tr(\Frob_{\mathfrak p}\mid V_4)=v(N\mathfrak p)$.
There are explicit four-dimensional $\Q_2$-representations $W_4^+$ and
$W_4^-$ of $G_\Q$, unramified outside $2,3$, whose restrictions to $G_K$
after scalar extension are $V_4$, with
\[
 W_4^-\simeq W_4^+\otimes\chi_{-3}.
\]
Both have a perfect alternating pairing with values in $\Q_2(-5)$
and are pure of weight five at every rational prime $p\ge5$.
Each is the $2$-adic realization of a rational Chow motive cut out on a
smooth projective fivefold, with Hodge type $(5,0)^2+(0,5)^2$;
see Theorem~\ref{thm:scholl-realization}.
For $p\ge5$, $p\equiv2\pmod3$, choose $d\in\F_p$ nonsquare,
$r^2=d$ in $\F_{p^2}$, and a primitive cubic character $\rho$ there.
Put
\[
 E'_x:Y^2=X(X^2-2xX+x^2-d),\quad a'_p(x)=p+1-\#E'_x(\F_p),
\]
\[
 S_p=\sum_{x\in\F_p}\rho(x+r)
       \bigl(a'_p(x)^4-3p\,a'_p(x)^2+p^2\bigr).
\]
Then $S_p\in\Z$ is independent of these choices and
\[
 \Tr(\Frob_p\mid W_4^+)=-S_p-1,\qquad
 \Tr(\Frob_p\mid W_4^-)=S_p+1.
\]
For $W_4^-$, set $c_1(p)=v(p)$ at split primes and $c_1(p)=S_p+1$
at inert primes, and set $c_2(p)=(c_1(p)^2-v(p^2))/2$ with the
character convention above. Then
\[
 \det(1-\Frob_pT\mid W_4^-)
 =1-c_1(p)T+c_2(p)T^2-p^5c_1(p)T^3+p^{10}T^4.
\]
\end{theorem}
\begin{theorem}[Local data and strong irreducibility]\label{thm:main-local}
Let $W_{4,\ell}^{\pm}$ be the realizations defined by the same rational
projectors as $W_4^{\pm}$ in Theorem~\ref{thm:scholl-realization}.
At $2$, for every $\ell\ne2$, both have zero inertia invariants and
Swan conductor zero. At $3$, for every $\ell\ne3$, the plus descent
has a one-dimensional invariant space with geometric Frobenius $3$,
whereas the minus descent has no invariants; both Swan conductors vanish.
Consequently the local Euler polynomials and Artin exponents are
\[
\begin{array}{c|cc|cc|c}
 &P_2(T)&a_2&P_3(T)&a_3&N\\\hline
W_4^+&1&4&1-3T&3&432\\
W_4^-&1&4&1&4&1296
\end{array}
\]
Here $N$ is the product of the local conductor powers, using a coefficient
prime different from the residue characteristic for each local calculation.
For every $\ell\ne2,3$, it is the prime-to-$\ell$ Artin conductor of
$W_{4,\ell}^{\pm}$. The displayed local data are independent of $\ell$
within their stated ranges.
For every $\ell\ne3$, both realizations are absolutely irreducible and
remain so after restriction to any finite number-field extension of $\Q$.
\end{theorem}
The local assertions are proved in Sections~\ref{sec:local-three} and
\ref{sec:local-two}, the conductor deduction in Section~\ref{sec:conductors},
and strong irreducibility in Section~\ref{sec:irreducibility}.
The last assertion includes the original $2$-adic realizations; no
irreducibility claim at coefficient prime $3$ is made here.

The construction, pairing and inert calculation are proved in
Sections~\ref{sec:geometry}, \ref{sec:descent}, \ref{sec:scholl} and \ref{sec:inert}.
The cubic second-moment comparison is treated separately in
Section~\ref{sec:cubic}; its modular conclusion uses the explicit exact
finite checks stated and independently recomputed there. These rank-two and rank-four objects
must not be conflated.

\subsection{Scope of this revision}
This consolidated research draft continues the version-2 preprint
\cite{WalkeV2}, with its earlier predecessor \cite{WalkeV1} retained for
provenance. It brings the local-at-$3$, dyadic tameness, dyadic invariant,
and strong-irreducibility arguments into the same manuscript.
The earlier blanket open obligations for these local data are superseded
by Theorem~\ref{thm:main-local}. Section~\ref{sec:claims} records the
current proof dependencies and unresolved questions. This draft has
internal AI-assisted review; external specialist review remains outstanding.
It does not certify the archived floating-point tables or promote their
functional-equation sign and special values to theorems.

The supplementary corrections retain known trace-formula consequences,
the Jacobi-sum normalization, and an analytic period substitution. Those
statements are explicitly distinguished from the main construction and
from numerical special-value observations. No classical elliptic BSD rank
or leading-term theorem is asserted.
